\clearpage
\fancyhead[L]{\small NEPTUN.}
\hypertarget{native-477-lower}{}
\begin{center}
\rule{0.55\linewidth}{0.4pt}\\[2pt]
Astronomische Nachrichten. Band XXVI. Nr. 610. Seite 153. 154. 1847 August 30.\\[2pt]
\rule{0.55\linewidth}{0.4pt}
\end{center}
\begin{flushright}Göttingen, 1847 Aug. 15.\end{flushright}

Beobachtungen des Neptun am Meridiankreise von \textsc{Gauss}.

\begin{center}
\setlength{\tabcolsep}{3pt}
\begin{tabular}{l|r|r|r}
 & Mittlere Zeit & Gerade Aufsteig. & Abweichung \\\hline
1847 Aug. 12 & $12^{\rm u}\ 43^{\rm m}\ 5^{\rm s}9$ & $331^\circ\ 34'\ 57''1$ & $-12^\circ\ 18'\ 12''1$ \\
\hfill 13 & $12\ 39\ 3.8$ & $331\ 33\ 27.0$ & $-12\ 18\ 44.6$ \\
\hfill 14 & $12\ 35\ 1.3$ & $331\ 31\ 50.1$ & $-12\ 19\ 19.5$ \\
\end{tabular}
\end{center}

\clearpage
\fancyhead[L]{\small IRIS.}
\SourceLeaf{478}
\begin{center}
\rule{0.55\linewidth}{0.4pt}\\[2pt]
Nachrichten der Königl. Gesellsch. d. Wiss. zu Göttingen. Nr. 11. 1847 September 13.\\[2pt]
\rule{0.55\linewidth}{0.4pt}
\end{center}

Von der am 13. August durch Hrn. \textsc{Hind} in London gemachten Entdeckung eines neuen Planeten erhielt ich die erste Kenntniss aus einem Pariser Zeitungsblatt vom 16 August und beobachtete ihn in Folge dieser Mittheilung den 21. August im Meridian. Da jene Nachricht nur zwei um Eine Stunde von einander abliegende Positionen enthalten hatte, so blieb anfangs einige Ungewissheit wegen der Identität, die jedoch bald nachher nach Eingang weiterer englischer Beobachtungen vom 14. und 15 August beseitigt wurde. Seitdem ist der Planet, welchem bekanntlich der Name Iris beigelegt ist, an jedem Abend, wo die Witterung es erlaubt hat, von mir am Meridiankreise beobachtet. Folgende Positionen habe ich bisher erhalten.

\begin{center}
\setlength{\tabcolsep}{4pt}
\begin{tabular}{lr|r|r|r}
\multicolumn{2}{c|}{1847} & \shortstack{Mittlere Zeit\\in Göttingen} & \shortstack{Gerade\\Aufsteigung} & \shortstack{Südliche\\Abweichung} \\\hline
August & 21 & $9^{\rm u}\ 52^{\rm m}\ 54^{\rm s}1$ & $297^\circ\ 47'\ 20''8$ & $13^\circ\ 42'\ 19''9$ \\
 & 28 & $9\ 21\ 16.6$ & $296\ 45\ 38.2$ & $13\ 54\ 47.9$ \\
 & 29 & $9\ 16\ 53.3$ & $296\ 38\ 46.6$ & $13\ 56\ 34.2$ \\
Sept. & 4 & $8\ 51\ 12.0$ & $296\ 7\ 14.1$ & $14\ 6\ 3.8$ \\
 & 6 & $8\ 42\ 54.5$ & $296\ 0\ 46.5$ & $14\ 8\ 56.1$ \\
\end{tabular}
\end{center}

Der Planet erschien in diesen Beobachtungen (bei meistens nicht besonders günstiger Luft) nur mit der Helligkeit eines Sternes 9ter Grösse.

Gleich nachdem die Beobachtung vom 28. August reducirt war, unternahm Hr. Prof. \textsc{Goldschmidt} die Berechnung der elliptischen Elemente, wobei er meine beiden Beobachtungen vom 21. und 28. August und eine in Cambridge von Hrn \textsc{Challis} gemachte vom 14. August zum Grunde legte. Folgendes sind seine Resultate.

\begin{center}
\small
\setlength{\tabcolsep}{2.5pt}
\begin{tabular}{llr}
\multicolumn{2}{l}{\shortstack[l]{Mittlere Anomalie 1847 August 14. $0^{\rm u}$,\\Mittl. Göttinger Zeit}} & $288^\circ\ 53'\ 21''9$ \\
Länge des Perihels & \mbox{.\ .\ .\ .\ .\ .\ .\ .\ .\ .\ .\ .\ .\ .} & $44\ 6\ 49.3$ \\
Länge des aufsteigenden Knoten & \mbox{.\ .\ .\ .\ .\ .\ .\ .\ .} & $260\ 27\ 58.8$ \\
\multicolumn{3}{l}{beide vom scheinbaren Aequinoctium des 14. August an gerechnet} \\
Mittlere tägliche siderische Bewegung & \mbox{.\ .\ .\ .\ .\ .\ .} & $982''5907$ \\
Logarithm der halben grossen Axe & \mbox{.\ .\ .\ .\ .\ .\ .\ .} & $0.3717560$ \\
Excentricität & \mbox{.\ .\ .\ .\ .\ .\ .\ .\ .\ .\ .\ .\ .\ .\ .\ .} & $0.2135543$ \\
Neigung der Bahn & \mbox{.\ .\ .\ .\ .\ .\ .\ .\ .\ .\ .\ .\ .\ .} & $5^\circ\ 29'\ 55''8$ \\
\end{tabular}
\end{center}

Hr. \textsc{Gould} aus Boston, der sich bei uns mit grossem Eifer den astronomischen Studien widmet, hat gleichfalls aus denselben Beobachtungen eine elliptische Bahn berechnet, die von der obigen wenig verschieden ist. Auch stimmt diese mit den spätern bis jetzt gemachten Beobachtungen noch sehr gut überein.

Zweierlei ist hieraus bereits zu schliessen. Erstens ist die planetarische Natur der Iris, und dass ihre Bahn, eben so wie die der Ceres, Pallas, Juno, Vesta, Asträa und Hebe, zwischen den Bahnen des Mars und des Jupiter liegt, zur Gewissheit gebracht. Zweitens wird voraussichtlich der Planet bei seiner Erscheinung im Jahr 1848 bedeutend heller sein, als im gegenwärtigen.

\noindent Göttingen, den 7. September 1847.\hfill\textsc{Gauss}.
